\chapter*{ЗАКЛЮЧЕНИЕ}
\addcontentsline{toc}{chapter}{ЗАКЛЮЧЕНИЕ}
\label{ch:conclusion}

В ходе выполнения выпускной квалификационной работы поставленные задачи были решены в полном объеме, а цель исследования достигнута. Был проведен комплексный теоретико-численный анализ data-driven методов восстановления динамики (DMD/DMDc, SINDy/SINDy-c и EDMD).

Основные результаты работы:
\begin{enumerate}
    \item Разработана модульная программная архитектура (\texttt{dynsys-methods-evaluator}), предоставляющая формализованный протокол для унифицированного тестирования алгоритмов на синтетических наборах данных с контролируемым уровнем шума. Данный пайплайн обладает высокой расширяемостью для интеграции новых методов.
    \item Сформирована эмпирическая карта применимости алгоритмов. Показано, что метод DMDc оптимален исключительно для линейных систем. Метод EDMD позволяет моделировать нелинейную динамику, но крайне чувствителен к уровню измерительного шума из-за мультипликативного эффекта при нелинейном расширении признаков. 
    \item Установлено, что алгоритм SINDy является наиболее робастным решением для идентификации систем со сложной топологией (включая хаотический аттрактор Лоренца) при наличии шума, при условии обеспечения достаточного объема обучающей выборки.
\end{enumerate}

Сформированные материалы воспроизводимости, включая скрипты генерации, модули оценки и графики, выложены в открытый доступ, что позволяет использовать результаты работы для выбора оптимального алгоритма в реальных прикладных задачах без необходимости физического моделирования.